%%-*-latex-*-

\vspace*{\stretch{1}}
\begin{figure}[H]
\centering
\fbox{\includegraphics[scale=0.4]{amitabha.eps}}
\scaption{Le Vainqueur Amit\a={a}bha}
\end{figure}
\vspace*{\stretch{1}}

\newpage

\titlesec{Transfert de la conscience}
\addcontentsline{toc}{section}{\protect\numberline{}Transfert de la conscience}

\bigskip\bigskip

\begin{tabbing}
\=\stib{mgon po 'od dpag med la gsol ba 'debs\notsheg:}\\
\>\lit{g�npo} \quad \=\lit{�}---------\=\lit{pagme} \=\lit{la}
\=\lit{s�lwa} \lit{deb}\\
\>\lex{protecteur} \>\lex{lumi�re} \>\lex{infinie} \>\lex{�}
\>\lex{je prie}\\
\>\trans{� Protecteur Amit\a={a}bha,\sendnote{Le Vainqueur (ou
  Protecteur) Amith\={a}bha est le repr�sentant de la famille du lotus
  (cf. note~5). Il sert ici de support � une pratique de
transfert de la conscience du mourant dans la Terre Pure de l'Ouest
(Sukh\={a}vat\={\i}), sur laquelle il r�gne. Il est
repr�sent� assis en lotus sur un tr�ne de lotus; ses mains font le
sceau (\emph{mudr\={a}}) du recueillement m�ditatif, et tiennent
un bol � aum�nes empli d'ambroisie. Il est flanqu�
de deux acolytes (\emph{bodhisattva}): � sa droite,
Avalokite\'{s}vara, blanc et tenant une fleur de lotus par la tige; �
sa gauche, Mah\={a}sth\={a}mapr\={a}pta (identifi� � Vajrap\={a}\d{n}i
d'apparence paisible dans le Vajray\={a}na), bleu et tenant un
\emph{vajra}. Ces deux assistants repr�sentent respectivement la
compassion en tant que sagesse et la compassion en tant qu'�nergie
salvifique. La rutilance illimit�e et isotrope d'Amit\={a}bha
caract�rise la compassion sans bornes et impartiale, pr�sente
potentiellement en chaque �tre (cf. l'\emph{Appel du ma�tre qui est au
  loin}, page~\pageref{compassion}) pour le b�n�fice d'autrui. Le
transfert de la conscience n'est ici qu'effleur� et n'est r�ellement mis
en pratique qu'apr�s que les pr�liminaires extraordinaires ont �t�
accomplis.} je vous prie}\\
\>\stib{zab lam 'pho ba 'byongs bar byin gyis rlobs\notsheg:}\\
\>\lit{sab}--------\=\lit{lam} \=\lit{phowa} \quad \=\lit{jong}-\lit{bra}
\=\lit{jin} \qquad \quad \=\lit{gyi} \=\lit{lob}\\
\>\lex{profonde} \>\lex{voie} \>\lex{tib. \emph{phowa}}
\>\lex{parfait dans} \>\lex{b�n�diction} \>\lex{par}
\>\lex{accordez}\\
\>\trans{de me b�nir afin que j'accomplisse la voie profonde du
transfert}\\
\>\trans{de conscience!}
\end{tabbing}

\rep{� r�citer trois fois.}

\newpage

\vspace*{\stretch{1}}
\begin{figure}[H]
\centering
\fbox{\includegraphics[scale=0.7]{samantabhadra.eps}}
\scaption{L'�veill� primordial Samantabhadra}
\end{figure}
\vspace*{\stretch{1}}
